\section{Cote 136: two verified parts of the weighted PD totalization attempt}
\label{sec:136twpartial}

\paragraph{Construction and scope.}
Let $B$ be the standard integral theory AQ-136-DP-MODEL and
$S=\Gamma_{\mathbb Z}(T)$. For an augmented cosimplicial
strict PD differential $S$-algebra $A^{-1}\to A^\bullet$,
define the weighted PD end
\[
 \operatorname{TW}_B(A)=
  \left\{(\omega_r)_{r\geq0}\in\prod_{r\geq0}
       (B_r\otimes_S A^r):
 (B(\alpha)\otimes1)\omega_s
       =(1\otimes A(\alpha))\omega_r\ \text{for }\alpha:[r]\to[s]
  \right\}.                                               \tag{136.TW1}
\]
Weights and form degrees are totalized in each tensor factor;
the differential is $d_B\otimes1+(-1)^{\deg_B}1\otimes d_A$.
The tensor has its usual strict PD structure: $B_r$ is a
free special-PD de Rham algebra over $S$ on $r$ variables.
The defining maps preserve that structure, and limits are
taken degreewise. No assertion that this end is a derived
end is implicit in its definition.

Take the cofree resolution $R=GU$ of
ARCH-136-ASSOCIATIVE-TERMINAL and define a simplicial PD
algebra
\[
       \mathsf T(C)_n=
       \operatorname{TW}_B((R^{\bullet+1}C)_n).
                                                               \tag{136.TW2}
\]
There are natural strict maps
$C\to\mathsf T(C)$ and
$\mathsf T(\underline S)\to\mathsf T(C)$.
If $\mathsf T(C)$ were admissible for every admissible $C$,
these would be actual PD common targets. The results below
prove the full vertical condition and the constant-input
case. The horizontal condition for general $C$ is not proved.
This is the original uniformity issue of cote 136, physical
pp.~117--121, footers 116--120, scans 117, 119 and 121.

\paragraph{ARCH-136-TW-VERTICAL: an integral contraction for split cosimplicial input.}
Suppose $A^{-1}\to A^\bullet$ is split as an augmented
cosimplicial differential $S$-module by extra codegeneracies
$t_m:A^m\to A^{m-1}$, $m\geq0$, commuting with the internal
differential. Then
\[
                A^{-1}\longrightarrow\operatorname{TW}_B(A)
                                                               \tag{136.TW3}
\]
is a weight-preserving chain-homotopy equivalence. In
particular $C_n^{w,*}\to\mathsf T(C)_n^{w,*}$ is such an
equivalence in every simplex and every weight. Thus the
vertical admissibility cutoffs of $\mathsf T(C)$ may be
taken equal to those of $C$, without any freeness assumption.

\begin{proof}
We first construct an ordered, integral contraction,
since the usual radial formula would introduce unnecessary
denominators. Regard vertex zero of $\Delta[r+1]$ as an apex.
Eliminating its coordinate gives
\[
 B_{r+1}=S\langle x_0,\ldots,x_r\rangle_{\rm PD}
                                \otimes\Lambda(dx_0,\ldots,dx_r).
\]
Here the displayed $x_i$ are the remaining, ordered base
vertices; the eliminated apex coordinate is $T-\sum_i x_i$.
For one variable let $\varepsilon_i$ set $x_i,dx_i$ to zero,
and let
\[
 h_i(x_i^{[a]}dx_i)=x_i^{[a+1]},\qquad h_i(x_i^{[a]})=0.
\]
The tensor contraction
\[
 h=\sum_{i=0}^r
       \varepsilon_0\otimes\cdots\otimes\varepsilon_{i-1}
                      \otimes h_i\otimes1\otimes\cdots\otimes1
                                                               \tag{136.TW4}
\]
satisfies $d_Bh+hd_B=1-\varepsilon_{\rm apex}$. The
ordinary tensor signs in a term with earlier augmentations
are trivial because those augmentations vanish on positive
form degree. The homotopy identity telescopes from the
one-variable identities.

The important naturality is for maps of cones induced
by order-preserving maps of their bases. Such a map
substitutes a sum of one consecutive block of variables
for each base variable; empty blocks are allowed.
Formula (136.TW4) is natural under these substitutions.
Here is a monomial check. On a block of variables, the
contraction is zero on a function without a differential.
On the divided-power one-form it gives
\[
 h\left(\gamma_a(\textstyle\sum_i x_i)
                          d(\textstyle\sum_i x_i)\right)
                   =\gamma_{a+1}(\textstyle\sum_i x_i).
                                                               \tag{136.TW5}
\]
Expand the divided-power sum. A summand involving $dx_i$
survives (136.TW4) only if all exponents before $i$ are
zero; it then raises the exponent of $x_i$ by one.
Every divided-power monomial of total exponent $a+1$
occurs exactly once, indexed by its first nonzero exponent.
This proves (136.TW5) integrally. In an ordered product
of blocks, the homotopy acts in the first block on which
the preceding augmentations do not already force zero.
The calculation just given and the augmentation identities
therefore prove naturality on every PD/exterior monomial.

Let
\[
 I_r=B(\delta^0)h:B_{r+1}\longrightarrow B_r[-1],
                                                               \tag{136.TW6}
\]
where $\delta^0:[r]\to[r+1]$ includes the base opposite
the apex. Then
\[
 d_BI_r+I_rd_B=B(\delta^0)-\varepsilon_{\rm apex},
 \qquad
 B(\alpha)I_s=I_r B(\operatorname{cone}\alpha).
                                                               \tag{136.TW7}
\]
In the first identity the apex evaluation is followed
by the coefficient inclusion into $B_r$. The second
is exactly the block-sum naturality proved above.

The extra-codegeneracy identities include
\[
 t_{r+1}A(\delta^0)=1,
 \qquad
 t_{s+1}A(\operatorname{cone}\alpha)
                       =A(\alpha)t_{r+1}.
                                                               \tag{136.TW8}
\]
Define a degree-minus-one linear map on the end by
\[
                   (H\omega)_r=(I_r\otimes t_{r+1})\omega_{r+1}.
                                                               \tag{136.TW9}
\]
The second identities in (136.TW7)--(136.TW8), together
with the end relation for $\operatorname{cone}\alpha$,
show that (136.TW9) again satisfies every end relation.
There are no infinite sums: the formula uses one component
of $\omega$ and a finite tensor contraction in that component.
It preserves the total weight.

Let $\pi(\omega)=t_0\omega_0$, using $B_0=S$, and let
$\iota:A^{-1}\to\operatorname{TW}_B(A)$ be the constant
coefficient family induced by the augmentation. The
internal $d_A$ terms cancel in $dH+Hd$, because $I_r$
lowers the $B$ degree by one and $t$ commutes with $d_A$.
The first term of (136.TW7) gives $\omega_r$, by the
end relation for $\delta^0$ and $t_{r+1}A(\delta^0)=1$.
For the apex term use the end relation for the vertex
$v_0:[0]\to[r+1]$. The extra-codegeneracy identities give
\[
         t_{r+1}A(v_0)=a_r t_0,
\]
where $a_r:A^{-1}\to A^r$ is the augmentation map in
that degree. Thus the apex term is exactly
$(\iota\pi\omega)_r$. We have proved
\[
                    dH+Hd=1-\iota\pi,
                    \qquad \pi\iota=1.
\]
This proves (136.TW3). For (136.TW2), evaluation in a
fixed original simplex makes the cofree resolution split
by the adjunction counit. Its extra codegeneracy is an
$S$-linear differential algebra map. The argument applies
in every weight without changing any bound. The resulting
$\pi$ and homotopy need not be natural in the original
simplex variable; only the inclusion is asserted to be
a simplicial map. The objectwise contraction is exactly
what the vertical admissibility condition requires.
\end{proof}

\paragraph{ARCH-136-TW-CONSTANT: the constant-input PD totalization is admissible.}
The simplicial PD algebra $\mathsf T(\underline S)$ is
admissible. Its vertical cohomology is $S$ in degree
zero in every weight and simplex. Its form-$q$ horizontal
row is acyclic whenever $w>q$, the standard cutoff for
$B$. There is a natural strict PD arrow
$B\to\mathsf T(\underline S)$.

\begin{proof}
In the coface convention $\delta^0=\eta_R$,
$\delta^1=R\eta$, the coordinates of
$(R^{r+1}\underline S)_n$ are indexed by strings
\[
 [l_r]\longrightarrow\cdots\longrightarrow[l_0]
                                      \longrightarrow[n].
\]
They form the $r$-simplices of
$K(n)=N((\Delta/\Delta[n])^{\rm op})$.
For example in degree one, $\delta^0$ selects the
composite simplex and $\delta^1$ the outer simplex,
which fixes this opposite-category convention.
The free $S$-basis of $B_r$, in each fixed form degree,
has only finitely many basis elements of weight at most
any specified bound. Thus its tensor product commutes
weightwise with the cofree product of copies of $S$.
The defining end is precisely
\[
               \mathsf T(\underline S)_n=B(K(n)).       \tag{136.TW10}
\]

More generally put $K(X)=N((\Delta/X)^{\rm op})$.
Its $r$-simplices are a disjoint union of copies of
$X_{l_0}$ indexed by strings
$[l_r]\to\cdots\to[l_0]$. Consequently $K$ preserves
colimits and monomorphisms of simplicial sets.
For a simplicial module $M$, define
\[
                  (E_KM)_n=
                  \operatorname{Hom}_{s\operatorname{Ab}}
                            (\mathbb Z[K(\Delta[n])],M).
\]
Colimit preservation gives the corresponding adjunction
on maps out of free simplicial modules on arbitrary
simplicial sets.

If $NM$ is acyclic, then $N(E_KM)$ is acyclic.
Indeed any map $\partial\Delta[m]\to E_KM$ corresponds
to a map $\mathbb Z[K(\partial\Delta[m])]\to M$.
The inclusion into $\mathbb Z[K(\Delta[m])]$ has free
relative normalized chains, because $K$ preserves
monomorphisms. Extend the map successively over their
free basis cells: at every step the prescribed boundary
is a cycle in $NM$ and has a primitive by acyclicity,
including in degree zero. Dold--Kan gives an extension
over $\Delta[m]$. For a normalized cycle of $E_KM$,
apply this to the boundary with that cycle on its last
face and zero on the others; a filler is a normalized
primitive. This proves all the claimed homology vanishing.

Apply this to $M=B_*^{w,q}$, which is acyclic for $w>q$
by AQ-136-DP-MODEL. Equation (136.TW10) proves horizontal
admissibility with that same cutoff. Vertical admissibility
in every weight follows from ARCH-136-TW-VERTICAL with
$C=\underline S$.

Finally there is a natural first-vertex map $K(X)\to X$.
For a string ending in $\sigma\in X_{l_0}$, take the
vertices $0,\alpha_1(0),\alpha_1\alpha_2(0),\ldots$
of $[l_0]$; they are weakly increasing. Restricting
$\sigma$ to this ordered sequence defines the map,
and composition verifies its face and degeneracy
identities. Contravariance of $B$ gives the stated
strict PD arrow. Thus this constant-input model is
already connected to the standard admissible theory.
\end{proof}

\paragraph{Horizontal question left by these two partial results.}
For a general admissible $C$, (136.TW2) has the required
strict algebraic structure and, by the first result,
exactly the correct vertical bounds. Its horizontal
rows are ends with varying coefficient modules $C_l$.
They are not the constant-coefficient rows in
(136.TW10). Limits of contractible cofree rows need
not be acyclic, and the first coface is not simply
the cofree functor applied to a map of graded modules.
No passage from the two verified assertions to full
admissibility of $\mathsf T(C)$ is made here.

The subsequent checked module ARCH-136-TW-TOR resolves
this question negatively over $\mathbb Z$: for $C=B$,
the zero-form row retains $\operatorname{Tor}_1^{\Gamma(T)}(k,k)$
in arbitrarily large prime-power weights.
ARCH-136-TW-TENSOR extends the identification to every
form row. The present vertical and constant-input
statements remain valid.
